% Document and references
\ifdefined\anonmode
  \newcommand*{\PurePy}{[Our lang]\xspace}
\else
  \newcommand*{\PurePy}{PurePy\xspace}
\fi
\ifdefined\anonmode
  \newcommand*{\issue}[1]{\##1}
\else
  \newcommand*{\issue}[1]{\href{https://github.com/pure-py/pure-py-spec/issues/#1}{\##1}}
\fi
\newcommand*{\versionstring}{}
\newcommand*{\version}[3]{%
  \gdef\versionstring{#1.#2.#3}%
}
\newcommand*{\figref}[1]{Figure~\ref{fig:#1}}
\newcommand*{\figrefTwo}[2]{Figures \ref{fig:#1} and \ref{fig:#2}}

% Figure and rule presentation
\definecolor{shadecolor}{rgb}{0.9,0.9,0.9}
\newcommand*{\shadebox}[1]{\fcolorbox{lightgray}{shadecolor}{\raisebox{0pt}[0.60\baselineskip][0.05\baselineskip]{#1}}}
\newcommand*{\ruleName}[1]{\textcolor{gray}{\textnormal{\textsf{#1}}}}

% General notation
\newcommand*{\set}[1]{\{#1\}}
\newcommand*{\dom}{\mathrm{dom}}
\newcommand*{\seq}[1]{\vec{#1}}
\newcommand*{\cons}{\cdot}
\newcommand*{\B}{\mathbb{B}}

% Terms: expressions and comprehension qualifiers
\newcommand*{\exCall}[2]{#1(#2)}
\newcommand*{\exConstr}[2]{#1(#2)}
\newcommand*{\exBinOp}[3]{#1 \mathbin{#2} #3}
\newcommand*{\exUnaryOp}[2]{#1\,#2}
\newcommand*{\exAttr}[2]{#1\mathord{.}#2}
\newcommand*{\exSubscript}[2]{#1[#2]}
\newcommand*{\exCondExpr}[3]{#1\;\pyinline{if}\;#2\;\pyinline{else}\;#3}
\newcommand*{\exList}[1]{[#1]}
\newcommand*{\exTuple}[1]{(#1)}
\newcommand*{\exDict}[1]{\{#1\}}
\newcommand*{\exListComp}[2]{[\,#1\;#2]}
\newcommand*{\exLambda}[2]{\pyinline{lambda}\;#1\mathord{:}\;#2}
\newcommand*{\qualGuard}[1]{\pyinline{if}\;#1}
\newcommand*{\qualGenerator}[2]{\pyinline{for}\;#1\;\pyinline{in}\;#2}

% Terms: statements, blocks, and module bodies
\newcommand*{\exPass}{\pyinline{pass}}
\newcommand*{\exAssign}[2]{#1 = #2}
\newcommand*{\exAssert}[1]{\pyinline{assert}\;#1}
\newcommand*{\exAssertMsg}[2]{\pyinline{assert}\;#1, #2}
\newcommand*{\exReturn}[1]{\pyinline{return}\;#1}
\newcommand*{\exReturnNone}{\pyinline{return}}
\newcommand*{\exIf}[3]{\pyinline{if}\;#1:\;#2\;#3}
\newcommand*{\exIfElse}[4]{\pyinline{if}\;#1:\;#2\;#3\;\pyinline{else}:\;#4}
\newcommand*{\exElif}[2]{\pyinline{elif}\;#1:\;#2}
\newcommand*{\exMatch}[2]{\pyinline{match}\;#1\mathord{:}\;#2}
\newcommand*{\exCase}[2]{\pyinline{case}\;#1\mathord{:}\;#2}
\newcommand*{\exDef}[3]{\pyinline{def}\;#1(#2):\;#3}
\newcommand*{\exDataclass}[2]{@\pyinline{dataclass}\;\pyinline{class}\;#1:\;#2}
\newcommand*{\exDataclassExtend}[3]{@\pyinline{dataclass}\;\pyinline{class}\;#1(#2):\;#3}
\newcommand*{\exField}[1]{#1:\;\pyinline{Any}}
\newcommand*{\exImport}[1]{\pyinline{import}\;#1}
\newcommand*{\exFromImport}[2]{\pyinline{from}\;#1\;\pyinline{import}\;#2}
\newcommand*{\exSeq}[2]{#1\;#2}

% Terms: patterns
\newcommand*{\patConstr}[2]{#1(#2)}
\newcommand*{\patWild}{\text{\pytext{\textunderscore}}}
\newcommand*{\patList}[1]{[#1]}
\newcommand*{\patTuple}[1]{(#1)}

% Values and context entries
\newcommand*{\bind}[2]{#1\mathord{:}\,#2}
\newcommand*{\statusDefAssigned}{\mathsf{tt}}
\newcommand*{\statusNotDefAssigned}{\mathsf{ff}}
\newcommand*{\modStub}[1]{\mathsf{Mod}(#1)}
\newcommand*{\modLoaded}[2]{\mathsf{Mod}(#1, #2)}
\newcommand*{\clsEntry}[4]{\mathsf{Cls}_{#1}(#2, #3, #4)}
\newcommand*{\exLamClosure}[4][]{\mathsf{Lam}_{#1}(#2, #3, #4)}
\newcommand*{\exDefClosure}[4][]{\mathsf{Def}_{#1}(#2, #3, #4)}
\newcommand*{\exObj}[3]{\mathsf{Obj}_{#1}(#2, #3)}

% Context and environment operations
\newcommand*{\override}{\cdot}
\newcommand*{\merge}{\oplus}
\newcommand*{\meet}{\wedge}
\newcommand*{\extend}{\lhd}
\newcommand*{\disjunion}{\uplus}

% Result types and results
\newcommand*{\Assigns}[1]{\mathsf{Assigns}\;#1}
\newcommand*{\Returns}{\mathsf{Returns}}
\newcommand{\assigns}[1]{\mathsf{assigns}\;#1}
\newcommand{\returns}[1]{\mathsf{returns}\;#1}

% Metafunctions
\newcommand*{\assignsS}{\mathrm{assigns}}
\newcommand*{\capturesS}{\mathrm{captures}}
\newcommand*{\capturesE}{\mathrm{captures}_{\mathsf{e}}}
\newcommand*{\fvS}{\mathrm{fv}}
\newcommand*{\fvE}{\mathrm{fv}_{\mathsf{e}}}
\newcommand*{\binds}{\mathrm{binds}}
\newcommand*{\ownFields}{\mathrm{own\text{-}fields}}
\newcommand*{\fields}{\mathrm{fields}}
\newcommand*{\ancestors}{\mathrm{ancestors}}
\newcommand*{\imports}{\mathrm{imports}}
\newcommand*{\deps}{\mathrm{deps}}
\newcommand*{\submodules}{\mathrm{submodules}}
\newcommand*{\nameOf}{\mathrm{name}}
\newcommand*{\unwrap}{\mathrm{unwrap}}
\newcommand*{\dispatch}{\mathrm{dispatch}}
\newcommand*{\env}{\mathrm{env}}
\newcommand*{\resolveClass}{\mathrm{resolve\text{-}class}}

% Orders, relations, and judgement symbols
\newcommand*{\prefixOf}{\preceq}
\newcommand*{\properPrefixOf}{\prec}
\newcommand*{\subsumed}{\sqsubseteq}
\newcommand*{\matches}{\rightsquigarrow}
\newcommand*{\notmatches}{\not\rightsquigarrow}
\newcommand*{\names}{\mathrel{\mathsf{names}}}
\newcommand*{\importsR}{\mathrel{\mathsf{imports}}}
\newcommand*{\loadsTo}{\mathrel{\mathsf{loads\mbox{-}to}}}

% Open term typing arrow
\newcommand{\carrow}[1]{\xrightarrow{\; #1\; }}

% Change markup (changes + changebar packages).
% No strikethrough: ulem's \sout is neither line-breakable nor metric-neutral.
\colorlet{deletedgray}{gray!70}
\setdeletedmarkup{\textcolor{deletedgray}{#1}}
% Capitalised variants add a margin changebar; use the lowercase forms inside
% figures or math, where changebar spans misbehave. \leavevmode commits a run-in
% heading in the default colour before the markup colour starts.
\newcommand{\Added}[1]{\leavevmode\protect\cbstart\added{#1}\protect\cbend\xspace}
\newcommand{\Deleted}[1]{\leavevmode\protect\cbstart\deleted{#1}\protect\cbend\xspace}
% The control space after #1 separates the new text from the deleted text shown
% after it (a plain space at the start of #2 is eaten by the changes package).
\newcommand{\Replaced}[2]{\leavevmode\protect\cbstart\replaced{#1\ }{#2}\protect\cbend\xspace}
% Deleted display material (diagrams, figures) that cannot live inside \Deleted:
% greyed out with a changebar in draft mode, discarded in final mode.
\newenvironment{DeletedBlock}{\par\cbstart\color{deletedgray}}{\cbend\par}
% Added display material (tables, figures) that cannot live inside \Added: alignment
% and rule material is fragile inside the changes package's argument.
\newenvironment{AddedBlock}{\par\cbstart\color{blue}}{\cbend\par}
% Final mode bypasses the changes package: its \IfIsEmpty test expands the
% replacement text, and any macro with #1 in its body then causes a runaway \iffalse.
\ifdefined\finalmode
  \renewcommand{\Added}[1]{#1}%
  \renewcommand{\Deleted}[1]{}%
  \renewcommand{\Replaced}[2]{#1}%
  \renewcommand{\added}[2][]{#2}%
  \renewcommand{\deleted}[2][]{}%
  \renewcommand{\replaced}[3][]{#2}%
  \renewenvironment{DeletedBlock}{\setbox0\vbox\bgroup}{\egroup}%
  \renewenvironment{AddedBlock}{\par}{\par}%
\fi
